\subsection{Proof of Lemma~\ref{lem:fqe-renyi-slack}}\label{sec:sup-renyi-slack}
Notation is that of Section~\ref{sec:fqe-renyi-above-one}: $p=1/6$, $s=1/2$,
$f_q(x)=(x^q-x)/(1-q)$, $\ell_q$ the secant of $f_q$ through $p,s$, and $\widehat h_q$ its cubic Hermite
interpolant there.


The proof rests on an elementary bound for the function $\Phi_q$ that appears in the fourth derivative.

\begin{lemma}[An elementary fourth-derivative bound; proved, RA-1]
\label{lem:fqe-slack-fourth-derivative}
Write $\ell_q(x)=\mu_qx+\nu_q$ and, for $1<q\le5/4$, put
\[
 A(q)=\frac{2(q+1)}{3-q}\left(\frac1{q-1}-\mu_q\right),\qquad
 D(q)=\frac{4(2q-1)(3-2q)}{(q-1)(3-q)(2-q)},
\]
and $\Phi_q(x)=2\nu_q+A(q)x-D(q)x^q$.
Define $\Phi_1$ by continuity. Then $\Phi_q(x)<141/100$ for
$1\le q\le5/4$ and $0\le x\le1$.
\end{lemma}
\begin{proof}
Set $r=q-1$, $p=1/6$, $s=1/2$, and $\delta=r(2-r)(1-r)$.
The function $\Phi_q$ is concave in $x$. Its derivative at $p$ is
$p^rN_p/\delta$, where
\[
 N_p=(2+r)(1-r)(3\cdot3^r-1)-4(1+r)(1-4r^2).
\]
Since $\ln3>1$, $3^r\ge1+r$, and substitution gives
$N_p\ge11r^2+13r^3\ge0$. Thus its maximum on $[0,1]$ is attained
in $[p,1]$.

We show that $\Phi_q(x)$ increases with $q$ on this latter interval.
Three endpoint derivatives suffice. Define
\[
 \begin{split}
 N_0&=6(1-r)3^r-6+5r+7r^2,\\
 N_1&=(2+r)(1-r)(3\cdot2^{-r}-6^{-r})-4(1+r)(1-4r^2),\\
 J_0&=(N_0'-(\ln6)N_0)\delta-N_0\delta',\\
 J_p&=(N_p'-(\ln6)N_p)\delta-N_p\delta',\qquad
 J_1=N_1'\delta-N_1\delta'.
 \end{split}
\]
Primes here mean $r$ derivatives. Direct substitution into $\Phi_q$ gives
\[
 \partial_q\Phi_q(p)=\frac{p^rJ_0}{3\delta^2},\qquad
 \partial_q\partial_x\Phi_q(p)=\frac{p^rJ_p}{\delta^2},\qquad
 \partial_q\partial_x\Phi_q(1)=\frac{J_1}{\delta^2}.
\]
All three $J$ are positive for $0<r\le1/4$. Here are coarse bounds
that permit this assertion to be checked with small rational numbers.
Write $J=\sum_{k=2}^6c_kr^k+R$; the constant and linear coefficients
are zero. The five coefficients have the following lower bounds,
and $|R|\le Mr^7$:
\[
\begin{array}{c|rrrrr|r}
 &c_2&c_3&c_4&c_5&c_6&M\\\hline
% RA-TAYLOR-BEGIN
J_0&4&-6&15&-14&0&3\\
J_p&27&-2&-33&52&-29&3\\
J_1&4&86&-75&2&0&14\\
% RA-TAYLOR-END
\end{array}
\]
For completeness, these bounds require only
$69/100<\ln2<7/10$ and $109/100<\ln3<11/10$.
Both follow from the first five terms of
$\ln u=2\sum_{j\ge0}z^{2j+1}/(2j+1)$, $z=(u-1)/(u+1)$,
with tail at most $2z^{11}/[11(1-z^2)]$.
Expand the displayed $N$ and $J$ through degree six and substitute these
two logarithm intervals to get the five columns. To check the remainder,
each $J$ is a polynomial of degree at most six plus terms
$Q(r)e^{\alpha r}$ with $\deg Q\le5$.
For $Q(r)=\sum_jQ_jr^j$, its tail after degree six has absolute value at most
\[
 r^7 e^{|\alpha|/4}\sum_j
       \frac{|Q_j|\,|\alpha|^{7-j}}{(7-j)!}.
\]
Use $3^{1/4}<4/3$, $6^{1/4}<8/5$ and $2^{1/4}<6/5$ for the
three possible exponential factors. The resulting sums are respectively
less than $3,3,14$, giving the last column. In particular the table yields
$J_0>2r^2$, $J_p>24r^2$ and $J_1>3r^2$ on $0<r\le1/4$;
for the last row use $86r-75r^2\ge0$.

Now let $W(x)=\partial_q\partial_x\Phi_q(x)$ and
$k(q)=q(q-1)D(q)>0$. Differentiating $\partial_x^2\Phi_q=-k(q)x^{q-2}$ gives
\[
 W'(x)=-k(q)x^{q-2}\bigl(\partial_q\ln k(q)+\ln x\bigr).
\]
The bracket increases with $x$, so $W$ has no interior minimum.
Its endpoint values at $p,1$ are positive by the table. Thus $W\ge0$,
and $\partial_q\Phi_q(x)\ge\partial_q\Phi_q(p)>0$ on $[p,1]$.
It remains only to bound $\Phi_{5/4}$.

Put $a=6^{-1/4}$, $b=2^{-1/4}$, and $t=x^{1/4}$. Then
\[
 \Phi_{5/4}(x)=2(b-a)+\frac{36}{7}(3b-a)t^4-\frac{64}{7}t^5
 \le2(b-a)+\frac{36\cdot9^4}{35\cdot20^4}(3b-a)^5,
\]
by maximising the last two terms over $t\ge0$; their maximum is at
$t=9(3b-a)/20$. Integer fourth powers give $a>319/500$ and $b<841/1000$.
Consequently
\[
 \Phi_{5/4}(x)
 <\frac{203}{500}+\frac{36\cdot9^4}{35\cdot20^4}
                         \left(\frac{377}{200}\right)^5
 <\frac{141}{100}.
\]
Concavity in $x$ and monotonicity in $q$ give the assertion everywhere;
the endpoint $q=1$ follows by continuity.
\end{proof}

\begin{proof}[Proof of Lemma~\ref{lem:fqe-renyi-slack}]
Let $g=f_q-\ell_q$ and
$G=\widehat h_q-f_q-(100/141)g^2$. Both $G$ and $G'$ vanish at $p,s$.
Since $f_q^{(4)}=-q(2-q)(3-q)x^{q-4}<0$, differentiation gives
\[
 G^{(4)}(x)=-f_q^{(4)}(x)\left(1-\frac{100}{141}\Phi_q(x)\right)>0
 \qquad(x>0).
\]
The cubic Hermite remainder, with repeated nodes $p,p,s,s$, gives
$G(x)=G^{(4)}(\xi)(x-p)^2(x-s)^2/24\ge0$ for $x>0$;
continuity gives the result at zero. No derivative at zero is needed.
Finally $\sum_i\ell_q(\lambda_i)=S_q^*$, since $\ell_q$ is affine and
$\sum_i\lambda_i=s+3p=1$. Cauchy--Schwarz gives the summed bound.
\end{proof}
